\documentclass[a4paper]{article}
\usepackage{amsfonts}
\usepackage{amsmath}
\usepackage{amssymb}
\usepackage{graphicx}     

\begin{document}
  
\noindent \large Auteur: Sadia Nazary \\
\noindent \large Studierichting: Informatica\\
\noindent \large Oefeningenreeks 6A nr. 10.4\\

\medskip

\normalsize

Laat de vier termen zijn: $x$, $xq$, $xq^2$, $xq^3$\\

$\left\{
\begin{array}{l}
x+xq = 28\\
xq^2+xq^3 = 175\\
\end{array}
\right.$\\

$\Rightarrow x(1+q)=28$\\

$\Rightarrow xq^2(1+q)=175$\\

$\Rightarrow \dfrac{xq^2(1+q)}{x(1+q)}=\dfrac{175}{28}$\\

$\Rightarrow q^2=\dfrac{25}{4}$\\

$\Rightarrow q=\dfrac{5}{2},-\dfrac{5}{2}$\\

Voor $q = \dfrac{5}{2}$:\\

$x\left(1+\dfrac{5}{2}\right)=28 \Rightarrow \dfrac{7}{2}x=28 \Rightarrow x=8$\\

Dus de vier termen zijn:\\

$(8,20,50,125)$\\

Voor $q = -\dfrac{5}{2}$:\\

$x\left(1-\dfrac{5}{2}\right)=28 \Rightarrow -\dfrac{3}{2}x=28 \Rightarrow x=-\dfrac{56}{3}$\\

Dus de vier termen zijn:\\

$\left(-\dfrac{56}{3},\dfrac{140}{3},-\dfrac{350}{3},\dfrac{875}{3}\right)$\\

\begin{thebibliography}{999}
%\bibitem{a} Referentie
\end{thebibliography}
\end{document}